\documentclass[11pt]{article}

\usepackage[a4paper,margin=1in]{geometry}
\usepackage{amsmath,amssymb,bm}
\usepackage[T1]{fontenc}
\usepackage{lmodern}
\usepackage{microtype}
\usepackage{hyperref}

\hypersetup{
  colorlinks=true,
  linkcolor=black,
  urlcolor=blue,
  citecolor=black
}

\newcommand{\dd}{\mathrm{d}}
\newcommand{\Tr}{\operatorname{Tr}}
\newcommand{\Oh}{\mathcal{O}}
\newcommand{\ii}{\mathrm{i}}

\title{Tree Plus One-Loop Equations for the Quartic Oscillator}
\author{}
\date{}

\begin{document}
\maketitle

\section{Euclidean one-loop action}

Consider the Euclidean quartic oscillator
\begin{equation}
S_E[q]=\int \dd \tau\left[
\frac{m}{2}\dot q^{2}+\frac{1}{2}m\omega^2q^2+\frac{\lambda}{24}q^4
\right].
\end{equation}
Split
\begin{equation}
q(\tau)=Q(\tau)+\eta(\tau).
\end{equation}
Expanding around an arbitrary background, and assuming boundary terms vanish, gives
\begin{equation}
S_E[Q+\eta]
=S_E[Q]
+\int \dd\tau\,\eta(\tau)E_Q(\tau)
+\frac{1}{2}\int \dd\tau\,\eta\,\Oh_Q\,\eta+
\cdots,
\end{equation}
where
\begin{equation}
E_Q=-m\ddot Q+m\omega^2 Q+\frac{\lambda}{6}Q^3,
\end{equation}
and the quadratic fluctuation operator is
\begin{equation}
\Oh_Q=-m\frac{\dd^2}{\dd\tau^2}+m\omega^2+\frac{\lambda}{2}Q^2(\tau)
=m\left[-\frac{\dd^2}{\dd\tau^2}+\Omega^2(\tau)\right],
\end{equation}
with
\begin{equation}
\Omega^2(\tau)=\omega^2+\frac{\lambda}{2m}Q^2(\tau).
\end{equation}
In the effective-action construction the linear term is cancelled by the source. It also vanishes if the background is chosen to obey the classical equation. The one-loop effective action is
\begin{equation}
\Gamma_E[Q]=S_E[Q]+\frac{\hbar}{2}\Tr\ln \Oh_Q+O(\hbar^2).
\end{equation}
The factor of $m$ in $\Oh_Q$ gives only a field-independent determinant and will be dropped below.

After subtracting the field-independent $Q=0$ determinant, the derivative expansion of the determinant is
\begin{equation}
\Delta\Gamma_E^{\rm 1-loop}[Q]
=\hbar\int \dd\tau\left[
\frac{1}{2}\bigl(\Omega-\omega\bigr)
+\frac{\dot\Omega^2}{16\Omega^3}
+O(\partial_\tau^4)
\right].
\label{eq:det-derivative-expansion}
\end{equation}
Since
\begin{equation}
\dot\Omega=\frac{\lambda Q\dot Q}{2m\Omega},
\end{equation}
we obtain the local slow-background action
\begin{equation}
\Gamma_E[Q]\simeq
\int \dd\tau\left[
\frac{m}{2}\dot Q^2
+V_{\rm eff}(Q)
+\hbar\,\frac{\lambda^2Q^2\dot Q^2}{64m^2\Omega^5}
\right],
\label{eq:local-action-euclidean}
\end{equation}
where
\begin{equation}
V_{\rm eff}(Q)=
\frac{1}{2}m\omega^2Q^2+\frac{\lambda}{24}Q^4
+\frac{\hbar}{2}\bigl[\Omega(Q)-\omega\bigr],
\qquad
\Omega(Q)=\sqrt{\omega^2+\frac{\lambda}{2m}Q^2}.
\end{equation}
This local action is valid in the adiabatic regime
\begin{equation}
\frac{|\dot\Omega|}{\Omega^2}\ll 1,
\qquad
\frac{|\ddot\Omega|}{\Omega^3}\ll 1,
\qquad \ldots
\end{equation}

\section{Euclidean equation of motion}

Write the local action as
\begin{equation}
\Gamma_E[Q]\simeq \int \dd\tau
\left[\frac{1}{2}Z(Q)\dot Q^2+V_{\rm eff}(Q)\right],
\end{equation}
with
\begin{equation}
Z(Q)=m+\hbar\frac{\lambda^2Q^2}{32m^2\Omega^5}.
\end{equation}
Varying gives
\begin{align}
\frac{\partial L_E}{\partial Q}
&=\frac{1}{2}Z'(Q)\dot Q^2+V'_{\rm eff}(Q),\\
\frac{\partial L_E}{\partial \dot Q}
&=Z(Q)\dot Q,\\
\frac{\dd}{\dd\tau}\frac{\partial L_E}{\partial \dot Q}
&=Z(Q)\ddot Q+Z'(Q)\dot Q^2.
\end{align}
Therefore the Euclidean equation of motion is
\begin{equation}
\boxed{
-Z(Q)\ddot Q-\frac{1}{2}Z'(Q)\dot Q^2+V'_{\rm eff}(Q)=0
}
\label{eq:eom-euclidean-general}
\end{equation}
or equivalently
\begin{equation}
\boxed{
Z(Q)\ddot Q+\frac{1}{2}Z'(Q)\dot Q^2=V'_{\rm eff}(Q).
}
\end{equation}
The needed derivatives are
\begin{equation}
V'_{\rm eff}(Q)=m\omega^2Q+\frac{\lambda}{6}Q^3
+\hbar\frac{\lambda Q}{4m\Omega},
\end{equation}
\begin{equation}
Z'(Q)=\hbar\frac{\lambda^2}{32m^2}
\left[
\frac{2Q}{\Omega^5}-\frac{5\lambda Q^3}{2m\Omega^7}
\right].
\end{equation}
Thus the explicit slow-background tree plus one-loop Euclidean equation is
\begin{equation}
\boxed{
-\left[m+\hbar\frac{\lambda^2Q^2}{32m^2\Omega^5}\right]\ddot Q
-\frac{\hbar\lambda^2}{64m^2}
\left[
\frac{2Q}{\Omega^5}-\frac{5\lambda Q^3}{2m\Omega^7}
\right]\dot Q^2
+m\omega^2Q+\frac{\lambda}{6}Q^3
+\hbar\frac{\lambda Q}{4m\Omega}=0.
}
\label{eq:eom-euclidean-explicit}
\end{equation}
Here all quantities are evaluated at $Q=Q(\tau)$, and a dot denotes $\dd/\dd\tau$. The equation is accurate up to $O(\hbar^2)$ and up to omitted higher-derivative terms, schematically $O(\partial_\tau^4)$ in the action.

Before the slow-background expansion, the exact Euclidean one-loop equation is nonlocal:
\begin{equation}
\boxed{
-m\ddot Q+m\omega^2Q+\frac{\lambda}{6}Q^3
+\frac{\hbar\lambda}{2}Q(\tau)G_E(\tau,\tau)=0,
}
\end{equation}
where
\begin{equation}
\left[-m\frac{\dd^2}{\dd\tau^2}+m\omega^2+\frac{\lambda}{2}Q^2(\tau)\right]
G_E(\tau,\tau')=\delta(\tau-\tau').
\end{equation}
As a useful check, the diagonal Green function has the adiabatic expansion
\begin{equation}
G_E(\tau,\tau)
=\frac{1}{m}\left[
\frac{1}{2\Omega}
-\frac{\ddot\Omega}{8\Omega^4}
+\frac{3\dot\Omega^2}{16\Omega^5}
+O(\partial_\tau^4)
\right].
\end{equation}
Using $\dot\Omega=\lambda Q\dot Q/(2m\Omega)$, this gives
\begin{equation}
\frac{\hbar\lambda}{2}QG_E(\tau,\tau)
=
\hbar\frac{\lambda Q}{4m\Omega}
-\hbar\frac{\lambda^2Q^2}{32m^2\Omega^5}\ddot Q
-\frac{\hbar\lambda^2}{64m^2}
\left[
\frac{2Q}{\Omega^5}-\frac{5\lambda Q^3}{2m\Omega^7}
\right]\dot Q^2
+O(\partial_\tau^4).
\end{equation}
Thus the local equation \eqref{eq:eom-euclidean-explicit} is the derivative expansion of this nonlocal one-loop result.

\section{Real-time one-loop equations}

The real-time classical action is
\begin{equation}
S[Q]=\int \dd t\left[
\frac{m}{2}\dot Q^2
-\frac{1}{2}m\omega^2Q^2
-\frac{\lambda}{24}Q^4
\right],
\end{equation}
where now a dot denotes $\dd/\dd t$. The quadratic fluctuation operator is
\begin{equation}
D_Q=-m\frac{\dd^2}{\dd t^2}-m\omega^2-\frac{\lambda}{2}Q^2(t)+\ii 0.
\end{equation}
The in-out one-loop effective action is
\begin{equation}
\Gamma_{\rm in-out}[Q]
=S[Q]+\frac{\ii\hbar}{2}\Tr\ln D_Q+O(\hbar^2).
\end{equation}
Varying it gives
\begin{equation}
\boxed{
m\ddot Q+m\omega^2Q+\frac{\lambda}{6}Q^3
+\frac{\ii\hbar\lambda}{2}Q(t)G_{\rm F}(t,t)=0,
}
\label{eq:realtime-exact-inout}
\end{equation}
where
\begin{equation}
D_QG_{\rm F}(t,t')=\delta(t-t').
\end{equation}
Here $G_{\rm F}$ denotes the inverse of $D_Q$; with this convention the usual
Feynman two-point function is $\ii\hbar G_{\rm F}$.
For a constant background, $G_{\rm F}(t,t)=-\ii/(2m\Omega)$, and \eqref{eq:realtime-exact-inout} reduces to
\begin{equation}
m\ddot Q+m\omega^2Q+\frac{\lambda}{6}Q^3
+\hbar\frac{\lambda Q}{4m\Omega}=0.
\end{equation}

For real causal time evolution one should use the in-in, or Schwinger--Keldysh, effective action. At one loop its equation has the same tadpole form,
\begin{equation}
\boxed{
m\ddot Q+m\omega^2Q+\frac{\lambda}{6}Q^3
+\frac{\lambda}{2}Q(t)\,\langle \eta^2(t)\rangle_Q=0,
}
\end{equation}
where $\langle\eta^2(t)\rangle_Q$ includes the factor of $\hbar$ and is evaluated in the Gaussian theory with time-dependent frequency $\Omega^2(t)=\omega^2+\lambda Q^2(t)/(2m)$. For an adiabatic vacuum, $\langle\eta^2\rangle_Q=\hbar/(2m\Omega)+\cdots$.

In the same slow-background approximation used above, the real part of the local real-time one-loop action is obtained by Wick rotating \eqref{eq:local-action-euclidean}:
\begin{equation}
\Gamma_{\rm RT}[Q]\simeq
\int \dd t\left[
\frac{1}{2}Z(Q)\dot Q^2-V_{\rm eff}(Q)
\right],
\end{equation}
with the same functions $Z(Q)$, $V_{\rm eff}(Q)$ and $\Omega(Q)$ as above. Therefore the local real-time equation of motion is
\begin{equation}
\boxed{
Z(Q)\ddot Q+\frac{1}{2}Z'(Q)\dot Q^2+V'_{\rm eff}(Q)=0.
}
\label{eq:realtime-local-general}
\end{equation}
Explicitly,
\begin{equation}
\boxed{
\left[m+\hbar\frac{\lambda^2Q^2}{32m^2\Omega^5}\right]\ddot Q
+\frac{\hbar\lambda^2}{64m^2}
\left[
\frac{2Q}{\Omega^5}-\frac{5\lambda Q^3}{2m\Omega^7}
\right]\dot Q^2
+m\omega^2Q+\frac{\lambda}{6}Q^3
+\hbar\frac{\lambda Q}{4m\Omega}=0.
}
\label{eq:realtime-local-explicit}
\end{equation}
All quantities in \eqref{eq:realtime-local-explicit} are evaluated at $Q=Q(t)$. This equation is the real-time adiabatic, local form of the one-loop corrected equation. For strongly time-dependent backgrounds, the full real-time equation is nonlocal and, in the in-out formalism, can acquire an imaginary part describing nonadiabatic excitation of fluctuations.

\end{document}
